\documentclass[11pt,a4paper]{article}

\usepackage[utf8]{inputenc}
\usepackage[T1]{fontenc}
\usepackage{lmodern}
\usepackage[margin=1in]{geometry}
\usepackage{amsmath,amssymb,amsthm,mathtools}
\usepackage{booktabs}
\usepackage{array}
\usepackage{enumitem}
\usepackage{xcolor}
\usepackage{listings}
\usepackage{hyperref}

\hypersetup{
  colorlinks=true,
  linkcolor=blue!60!black,
  urlcolor=blue!60!black,
  citecolor=blue!60!black
}

\lstdefinelanguage{Lean4}{
  keywords={def, theorem, lemma, by, have, let, in, if, then, else, match, with, fun, intro, rcases, rw, exact, refine, apply, omega, decide, simp, dsimp, ext, subst, generalize, constructor, forall, exists, import, namespace, end, open, use, by_cases, split_ifs},
  keywordstyle=\color{blue!70!black}\bfseries,
  comment=[l]{--},
  morecomment=[s]{/-}{-/},
  commentstyle=\color{green!45!black}\itshape,
  basicstyle=\ttfamily\small,
  breaklines=true,
  columns=fullflexible,
  keepspaces=true,
  frame=single,
  rulecolor=\color{black!20},
  backgroundcolor=\color{black!2},
  literate=
    {ℕ}{{$\mathbb{N}$}}1
    {ℤ}{{$\mathbb{Z}$}}1
    {×}{{$\times$}}1
    {∀}{{$\forall$}}1
    {∃}{{$\exists$}}1
    {∈}{{$\in$}}1
    {≤}{{$\le$}}1
    {≥}{{$\ge$}}1
    {≠}{{$\ne$}}1
    {∧}{{$\land$}}1
    {∨}{{$\lor$}}1
    {→}{{$\to$}}1
    {⇒}{{$\Rightarrow$}}1
    {⟨}{{$\langle$}}1
    {⟩}{{$\rangle$}}1
    {₁}{{$_1$}}1
    {₂}{{$_2$}}1
    {·}{{$\cdot$}}1
}

\newtheorem{theorem}{Theorem}[section]
\newtheorem{lemma}[theorem]{Lemma}
\newtheorem{proposition}[theorem]{Proposition}
\newtheorem{definition}[theorem]{Definition}
\theoremstyle{remark}
\newtheorem{remark}[theorem]{Remark}

\setlength{\parindent}{0pt}

\title{\textbf{
  Cyclic Incremental Latin $n \times n$ Squares \\ 
  encoding \\
  Generalized NC-Coins Squares for $0 \le c \le n$
}}
\author{Marc Duponcheel and Luc Duponcheel
}

\date{\today}

\begin{document}

\maketitle

\begin{abstract}

\noindent For every natural number $n \notin \{0, 2, 3, 4, 6\}$, there exists a permutation 
$q : \mathrm{Fin}\,n \to \mathrm{Fin}\,n$ such that
$L_{i,j} = (i + q(j)) \bmod n$ is an \emph{Cyclic Incremental Latin Square}, encoding NC-Coins
Squares for $0 \le c \le n$.

\vspace{12pt}

An NC-Coins Square for $0 \le 1 \le n$ corresponds to a solution of $n$-Queens Problem.

\vspace{12pt}

\noindent We have an existence proof in {\tt Lean} and construction code in {\tt Haskell},
{\tt Julia} and {\tt Lean}.

\vspace{12pt}

\noindent See e.g. \href{https://mathoverflow.net/questions/510135/placing-n-cdot-m-coins-on-an-n-times-n-board-generalizing-the-n-queens-p}{MO\hspace{-0.5pt}\#\hspace{-0.5pt}510135}, 
\href{https://puzzling.stackexchange.com/questions/138939/place-as-many-queens-as-you-can}{PSE\hspace{-0.5pt}\#\hspace{-0.5pt}138939},
\href{https://mathoverflow.net/questions/516225/universal-board-encoding-placement-of-n-cdot-m-coins-on-an-n-times-n-board}{MO\hspace{-0.5pt}\#\hspace{-0.5pt}516225} and
\href{https://github.com/NofanTasi/incremental-nc-coins-problem}{GitHub}.

\vspace{12pt}

\noindent The proof uses the \emph{Injective-Majorization Principle} that, in combination with the
\emph{Pigeon Hole Principle}, replaces intensive checking inequalities on each diagonal to the
existence of a single function $f(i,j) \le L_{i,j}$ that is injective on every diagonal. 

\vspace{12pt}

\noindent The proof uses four families covering almost all natural number modulo~$6$ classes
\begin{itemize}
\item for $n \equiv 1,5 \pmod 6$ (with $n \ge 1$)
\item for $n \equiv 2,4 \pmod 6$ (with $n \ge 10$)
\item for $n \equiv 0 \pmod 6$ (with $n \ge 12$)
\item for $n \equiv 3 \pmod 6$ (with $n \ge 27$)
\end{itemize}

\vspace{12pt}

\noindent For all four families a Cyclic Incremental Latin Square can be constructed. Constructing
a Cyclic Incremental Latin Square for the four uncovered natural numbers, $n \in \{8, 9, 15, 21\}$,
is handled separately. Likewise, constructing a (non-Cyclic) Incremental Latin Square for the two
remaining natural numbers, $n \in \{4, 6\}$, is also handled separately.

\end{abstract}

\section{Cyclic Incremental Latin Squares}

\subsection{From NC-Coins Squares to Cyclic Incremental Latin Squares}

NC-Coins Squares are subsets $S_c \subseteq \{0, \dots, n-1\}^2$ of $nc$ cells of an $n \times n$
square such that every row, column, sum-diagonal ($i + j = s$), and difference-diagonal
($i - j = d$) contains at most $c$ cells of $S_c$.

\vspace{12pt}

Clearly, for $c = 1$, they correspond to solutions of $n$-Queens Problems.

\vspace{12pt}

The construction of Le, Li and Wang~\cite{LLW1990} constructs $S_c$'s separately for each  $c$.

\vspace{12 pt}

\subsection{The Conjecture}

Every permutation $q : \mathrm{Fin}\,n \to \mathrm{Fin}\,n$ defines a \emph{cyclic Latin square} (a column-permuted Cayley table of $\mathbb{Z}/n\mathbb{Z}$):
\[
  L_{i,j} = (i + q(j)) \bmod n \qquad (0 \le i, j < n).
\]
Because each column is a cyclic shift of $(0, 1, \dots, n-1)$ by $q(j)$ and $q$ is a bijection, every row and every column of $L$ is automatically a permutation of $\{0, \dots, n-1\}$.

\vspace{12 pt}

The conjecture is that for all $n \ge 1$ except $n \in \{2, 3, 4, 6\}$, there exists a column permutation $q$ such that $L_{i,j} = (i + q(j)) \bmod n$ is an incremental Latin square. (For $n \in \{2, 3\}$, even the $k = 0$ layer---the classical $n$-queens problem---has no solution; for $n \in \{4, 6\}$, exhaustive finite search confirms that no cyclic Latin square satisfies the diagonal bounds.)

\begin{theorem}[Cyclic Incremental Latin Square Conjecture]\label{thm:main}
For every integer $n > 0$ with $n \notin \{2, 3, 4, 6\}$, there exists a bijection $q : \mathrm{Fin}\,n \to \mathrm{Fin}\,n$ such that the cyclic Latin square
\[
  L_{i,j} = (i + q(j)) \bmod n
\]
satisfies:
\begin{enumerate}[nosep,label=(\roman*)]
  \item For every $i \in \mathrm{Fin}\,n$, the row map $j \mapsto L_{i,j}$ is bijective.
  \item For every $j \in \mathrm{Fin}\,n$, the column map $i \mapsto L_{i,j}$ is bijective.
  \item For every $0 \le k < n$ and every $s \in \mathbb{N}$,
  \[
    |\{(i,j) \in \mathrm{Fin}\,n \times \mathrm{Fin}\,n \mid i + j = s \;\land\; L_{i,j} \le k\}| \le k + 1.
  \]
  \item For every $0 \le k < n$ and every $d \in \mathbb{Z}$,
  \[
    |\{(i,j) \in \mathrm{Fin}\,n \times \mathrm{Fin}\,n \mid i - j = d \;\land\; L_{i,j} \le k\}| \le k + 1.
  \]
\end{enumerate}
\end{theorem}


\begin{thebibliography}{9}
\bibitem{LLW1990}
Mao-Hua Le, Wen Li, and Edward T. Wang,
\emph{A generalization of the $n$-queen problem},
Journal of Systems Science and Mathematical Sciences \textbf{3} (1990), no.~2, 183--192,
\url{https://sysmath.cjoe.ac.cn/jssc/CN/Y1990/V3/I2/183}.
\end{thebibliography}

\end{document}